% ── Collision-Inclusive Trajectory Optimization ─────────────

\begin{frame}[plain]
  \vfill
  \begin{center}
    {\LARGE\color{slate900}\textbf{Collision-Inclusive Trajectory Optimization}}\\[8pt]
    {\color{slate700}Planning for Failure}
  \end{center}
  \vfill
\end{frame}

% ── Motivation: Free-Flying Spacecraft ───────────────────
\begin{frame}{Application: Free-Flying Robotic Spacecraft}
  \small
  \begin{columns}[c]
    \begin{column}{0.45\textwidth}
      Small robotic spacecraft are used for automated {\color{blue600}monitoring and assistance}.

      \vskip6pt
      \textbf{Problem:} Traditional planners avoid all collisions.
      \vskip2pt
      But sometimes collisions are \textbf{unavoidable}.

      \vskip6pt
      \textbf{Goal:} Find the best strategy when collisions cannot be avoided.

      \vskip2pt
      \textbf{Secondary goal:} Achieve better performance by \textit{using} collisions.

      \vskip6pt
      \textbf{Approach:} Optimal trajectory planning with collision modes encoded as constraints.
    \end{column}
    \begin{column}{0.52\textwidth}
      \begin{center}
        % Placeholder: AERCam Sprint, Astrobee, Spheres photos
        {\color{slate400}\framebox(220,100){AERCam Sprint / Astrobee / Spheres}}

        \vskip6pt
        % Placeholder: Stanford free-flier lab
        {\color{slate400}\framebox(220,80){Stanford Free-flier Lab}}

        \vskip2pt
        {\scriptsize\color{slate400} Stanford's Free-flier Lab (credit: Marco Pavone, Andrew Bylard)}
      \end{center}
    \end{column}
  \end{columns}
  \source{Mote, Egerstedt, Feron, Bylard, Pavone --- ``Collision-Inclusive Trajectory Optimization,'' \textit{JGCD}, 2020}
\end{frame}

% ── Collision Model and Solution ─────────────────────────
\begin{frame}{Collision Model and Solution Framework}
  \small
  \begin{columns}[c]
    \begin{column}{0.50\textwidth}
      \textbf{Nominal Dynamics:}\quad $\ddot{s}_x = u_x$,\quad $\ddot{s}_y = u_y$,\quad $\ddot{\theta} = u_\theta$

      \vskip4pt
      \textbf{Collision Model} (empirical, algebraic):
      \vskip-6pt
      \footnotesize
      \[
        \begin{bmatrix} v_N^{i+1} \\ v_T^{i+1} \\ \omega^{i+1} \end{bmatrix}
        = \begin{bmatrix} a_{11} & 0 \\ 0 & a_{22} \\ 0 & a_{32} \end{bmatrix}
        \begin{bmatrix} v_N^i \\ v_{\text{rel}}^i \end{bmatrix}
      \]
      \vskip0pt
      \small
      Developed empirically: perform collisions, regression on data.

      \vskip6pt
      \textbf{Solution:} Mixed Integer Quadratic Program (MIQP)
      \vskip2pt
      \footnotesize
      \begin{itemize}
        \item Integer variables encode collision modes
        \item Global optimality, completeness
        \item LQR regulation to ideal trajectory
      \end{itemize}
    \end{column}
    \begin{column}{0.47\textwidth}
      \begin{center}
        % Placeholder: collision model regression plots
        {\color{slate400}\framebox(190,90){Collision model regression plots}}

        \vskip4pt
        % Placeholder: free-flyer schematic with thrusters
        {\color{slate400}\framebox(190,90){Free-flyer schematic}}

        \vskip2pt
        {\scriptsize\color{slate400} 8 thrusters, reaction wheel, CO$_2$ tanks}
      \end{center}
    \end{column}
  \end{columns}
  \source{Mote, Afman, Feron --- ``Robotic trajectory planning through collisional interaction,'' CDC 2017}
\end{frame}

% ── Safety through Collision-Inclusive Planning ──────────
\begin{frame}{Safety through Collision-Inclusive Planning}
  \small
  \textbf{1D example:}\quad $\ddot{x} = u$, \quad $u \in [-0.5,\, 0.5]$, \quad wall at $x = 0$

  \vskip6pt
  \begin{columns}[c]
    \begin{column}{0.30\textwidth}
      \begin{center}
        % Placeholder: reachable without collision
        {\color{slate400}\framebox(115,120){Reachable\\w/o collision}}

        \vskip2pt
        {\scriptsize\color{blue600} Reachable without collision}
      \end{center}
    \end{column}
    \begin{column}{0.30\textwidth}
      \begin{center}
        % Placeholder: feasible under MPC
        {\color{slate400}\framebox(115,120){Feasible\\under MPC}}

        \vskip2pt
        {\scriptsize\color{blue600} Feasible under MPC}
      \end{center}
    \end{column}
    \begin{column}{0.30\textwidth}
      \begin{center}
        % Placeholder: feasible under CT-MPC
        {\color{slate400}\framebox(115,120){Feasible\\under CT-MPC}}

        \vskip2pt
        {\scriptsize\color{blue600} Feasible under CT-MPC}
      \end{center}
    \end{column}
  \end{columns}

  \vskip6pt
  \textbf{Key insight:} Collision-tolerant planning \textbf{expands the feasible set}. States that are infeasible under collision-free planning become feasible when collisions are allowed.
\end{frame}

% ── Conclusion: Collisions → Feasibility → Safety ───────
\begin{frame}{Planning for Failure --- Summary}
  \small
  \begin{columns}[c]
    \begin{column}{0.55\textwidth}
      \textbf{Collision-tolerant planners can:}
      \vskip4pt
      \begin{itemize}
        \item Optimally mitigate damage with respect to a predefined \textit{damage function}
        \item Expand the feasible set for trajectory planning
        \item Serve as better backup strategies for RTA systems
      \end{itemize}

      \vskip8pt
      \begin{center}
        {\color{blue600}\textbf{Collisions}}\\
        $\Downarrow$\\
        {\color{blue600}\textbf{Feasibility}}\\
        $\Downarrow$\\
        {\color{blue600}\textbf{Safety}}
      \end{center}
    \end{column}
    \begin{column}{0.42\textwidth}
      \begin{center}
        % Placeholder: optimal damage mitigation simulation
        {\color{slate400}\framebox(180,160){Optimal damage mitigation\\simulation frame}}

        \vskip2pt
        {\scriptsize\color{slate400} Optimal path through obstacles\\with damage minimization}
      \end{center}
    \end{column}
  \end{columns}
  \source{Mote, Bylard, Egerstedt, Pavone, Feron --- ``Collision-Inclusive Trajectory Optimization,'' \textit{JGCD}, 2020}
\end{frame}
